\documentclass[aspectratio=169]{beamer}

\usetheme{Madrid}
\usecolortheme{default}

\usepackage{amsmath,amssymb,amsfonts,bm}
\usepackage{physics}
\usepackage{mathtools}

\title{Stochastic Reconfiguration in Quantum Monte Carlo}
\author{Morten Hjorth-Jensen}
\date{FYS4411/9411 Spring 2026}

\begin{document}

\frame{\titlepage}

\begin{frame}{Content}
\begin{itemize}
\item imaginary-time evolution and variational projection,
\item tangent-space expansion of the trial state,
\item derivation of the linear system \(S \, \delta \bm{\alpha} = -\delta\tau \, \bm{g}\),
\item interpretation of the covariance matrix \(S\),
\item relation to the natural gradient.
\end{itemize}
\end{frame}

\begin{frame}{Setup and notation}
We consider a variational wave function
\[
\ket{\Psi(\bm{\alpha})},
\qquad
\bm{\alpha}=(\alpha_1,\alpha_2,\dots,\alpha_p),
\]
depending on \(p\) real parameters.

\medskip

The normalized state is
\[
\ket{\overline{\Psi}(\bm{\alpha})}
=
\frac{\ket{\Psi(\bm{\alpha})}}
{\sqrt{\braket{\Psi(\bm{\alpha})}{\Psi(\bm{\alpha})}}}.
\]

We wish to optimize the parameters such that the energy
\[
E(\bm{\alpha})
=
\frac{\bra{\Psi(\bm{\alpha})}H\ket{\Psi(\bm{\alpha})}}
{\braket{\Psi(\bm{\alpha})}{\Psi(\bm{\alpha})}}
\]
is minimized.
\end{frame}

\begin{frame}{Imaginary-time evolution as a route to the ground state}
The exact imaginary-time evolution of a state is
\[
\ket{\Psi(\tau)}
=
e^{-\tau H}\ket{\Psi(0)}.
\]

Expanding the initial state in energy eigenstates,
\[
\ket{\Psi(0)}=\sum_n c_n \ket{n},
\qquad
H\ket{n}=E_n\ket{n},
\]
gives
\[
\ket{\Psi(\tau)}
=
\sum_n c_n e^{-\tau E_n}\ket{n}.
\]

If \(c_0\neq 0\), then for large \(\tau\),
\[
\ket{\Psi(\tau)} \propto e^{-\tau E_0}\ket{0},
\]
so imaginary-time propagation filters out the ground state.
\end{frame}

\begin{frame}{Why a variational projection is needed}
In Variational Monte Carlo we do not evolve an arbitrary state in the full Hilbert space.
Instead, we restrict ourselves to the variational manifold
\[
\mathcal{M}
=
\left\{
\ket{\Psi(\bm{\alpha})} \; \middle| \; \bm{\alpha}\in\mathbb{R}^p
\right\}.
\]

The exact state
\[
(1-\delta\tau H)\ket{\Psi(\bm{\alpha})}
\]
will in general \emph{not} belong to \(\mathcal{M}\).

\medskip

Idea of stochastic reconfiguration:
\begin{itemize}
\item apply one infinitesimal imaginary-time step,
\item project the resulting state back onto the variational manifold,
\item interpret the projection as a parameter update \(\bm{\alpha}\to \bm{\alpha}+\delta\bm{\alpha}\).
\end{itemize}
\end{frame}

\begin{frame}{First-order imaginary-time step}
For a small time step \(\delta\tau\),
\[
e^{-\delta\tau H}
=
1-\delta\tau H + \mathcal{O}(\delta\tau^2).
\]

Thus the exact propagated state is approximated by
\[
\ket{\Psi_{\mathrm{exact}}'}
\approx
(1-\delta\tau H)\ket{\Psi(\bm{\alpha})}.
\]

A convenient first-order, norm-preserving form is
\[
\ket{\Psi_{\mathrm{exact}}'}
\approx
\left[1-\delta\tau\bigl(H-E(\bm{\alpha})\bigr)\right]\ket{\Psi(\bm{\alpha})}.
\]

Subtracting \(E(\bm{\alpha})\) removes the trivial component parallel to \(\ket{\Psi}\).
\end{frame}

\begin{frame}{Why subtract \(E(\bm{\alpha})\)?}
Let
\[
E=\frac{\bra{\Psi}H\ket{\Psi}}{\braket{\Psi}{\Psi}}.
\]

Then
\[
\bra{\Psi}(H-E)\ket{\Psi}=0.
\]

Hence the vector
\[
(H-E)\ket{\Psi}
\]
is orthogonal to \(\ket{\Psi}\) in the sense relevant for norm changes to first order.

\medskip

Therefore,
\[
\left[1-\delta\tau(H-E)\right]\ket{\Psi}
\]
moves the state in a direction tangent to projective Hilbert space, removing the trivial norm-rescaling component.
\end{frame}

\begin{frame}{Variational state after a parameter update}
We now expand the variational state at shifted parameters:
\[
\ket{\Psi(\bm{\alpha}+\delta\bm{\alpha})}
=
\ket{\Psi(\bm{\alpha})}
+
\sum_k \delta\alpha_k
\pdv{}{\alpha_k}\ket{\Psi(\bm{\alpha})}
+\mathcal{O}(\delta\alpha^2).
\]

Define the tangent vectors
\[
\ket{\Psi_k}
\equiv
\pdv{}{\alpha_k}\ket{\Psi(\bm{\alpha})}.
\]

Then
\[
\ket{\Psi(\bm{\alpha}+\delta\bm{\alpha})}
=
\ket{\Psi}
+
\sum_k \delta\alpha_k \ket{\Psi_k}
+\mathcal{O}(\delta\alpha^2).
\]
\end{frame}

\begin{frame}{Projection principle: minimize the distance}
We compare
\[
\ket{\Psi_{\mathrm{var}}'}
=
\ket{\Psi}
+
\sum_k \delta\alpha_k \ket{\Psi_k}
\]
with
\[
\ket{\Psi_{\mathrm{exact}}'}
=
\left[1-\delta\tau(H-E)\right]\ket{\Psi}.
\]

Define the mismatch vector
\[
\ket{\Delta}
=
\ket{\Psi_{\mathrm{var}}'}
-
\ket{\Psi_{\mathrm{exact}}'}.
\]

To first order,
\[
\ket{\Delta}
=
\sum_k \delta\alpha_k \ket{\Psi_k}
+
\delta\tau(H-E)\ket{\Psi}.
\]

The SR update is obtained by minimizing
\[
\mathcal{D}
=
\braket{\Delta}{\Delta}.
\]
\end{frame}

\begin{frame}{Expansion of the distance functional}
Using
\[
\ket{\Delta}
=
\sum_k \delta\alpha_k \ket{\Psi_k}
+
\delta\tau(H-E)\ket{\Psi},
\]
we obtain
\[
\mathcal{D}
=
\left(
\sum_i \delta\alpha_i \bra{\Psi_i}
+
\delta\tau \bra{\Psi}(H-E)
\right)
\left(
\sum_j \delta\alpha_j \ket{\Psi_j}
+
\delta\tau (H-E)\ket{\Psi}
\right).
\]

Expanding and keeping terms up to first order in \(\delta\tau\),
\[
\mathcal{D}
=
\sum_{ij}\delta\alpha_i \delta\alpha_j \braket{\Psi_i}{\Psi_j}
+
2\delta\tau \sum_i \delta\alpha_i \Re \bra{\Psi_i}(H-E)\ket{\Psi}
+\mathcal{O}(\delta\tau^2).
\]
\end{frame}

\begin{frame}{Stationarity condition}
Differentiate \(\mathcal{D}\) with respect to \(\delta\alpha_i\):
\[
\pdv{\mathcal{D}}{(\delta\alpha_i)}
=
2\sum_j \braket{\Psi_i}{\Psi_j}\,\delta\alpha_j
+
2\delta\tau \Re \bra{\Psi_i}(H-E)\ket{\Psi}.
\]

Setting this equal to zero gives
\[
\sum_j \braket{\Psi_i}{\Psi_j}\,\delta\alpha_j
=
-\delta\tau \Re \bra{\Psi_i}(H-E)\ket{\Psi}.
\]

This is the projected imaginary-time equation in the tangent basis \(\{\ket{\Psi_i}\}\).
\end{frame}

\begin{frame}{Normalized-state formulation}
To remove the redundant component parallel to \(\ket{\Psi}\), introduce
\[
\ket{\overline{\Psi}}=\frac{\ket{\Psi}}{\sqrt{\braket{\Psi}{\Psi}}}.
\]

Define normalized tangent vectors:
\[
\ket{\overline{\Psi}_k}
=
\pdv{}{\alpha_k}\ket{\overline{\Psi}}.
\]

Then the projected equation becomes
\[
\sum_j
\braket{\overline{\Psi}_i}{\overline{\Psi}_j}
\,\delta\alpha_j
=
-\delta\tau
\bra{\overline{\Psi}_i}(H-E)\ket{\overline{\Psi}}.
\]

This form leads directly to the SR metric matrix.
\end{frame}

\begin{frame}{Logarithmic derivatives in configuration space}
Let \(x\) denote a many-body configuration and write the wave function as \(\Psi(x)\).

Define the logarithmic derivative operators
\[
O_k(x)
=
\frac{1}{\Psi(x)}
\pdv{\Psi(x)}{\alpha_k}
=
\pdv{\ln \Psi(x)}{\alpha_k}.
\]

Then
\[
\pdv{\Psi(x)}{\alpha_k}=O_k(x)\Psi(x).
\]

For Monte Carlo sampling with probability
\[
P(x)=\frac{|\Psi(x)|^2}{\braket{\Psi}{\Psi}},
\]
expectation values are
\[
\langle A \rangle = \sum_x P(x)\, A(x)
\]
(or integrals in the continuous case).
\end{frame}

\begin{frame}{Deriving the metric matrix \(S_{ij}\)}
Using normalized states, one finds
\[
\braket{\overline{\Psi}_i}{\overline{\Psi}_j}
=
\langle O_i^\ast O_j\rangle
-
\langle O_i^\ast\rangle \langle O_j\rangle.
\]

For real wave functions and real parameters this becomes
\[
S_{ij}
=
\langle O_i O_j\rangle
-
\langle O_i\rangle \langle O_j\rangle.
\]

Thus \(S_{ij}\) is the covariance matrix of the logarithmic derivatives.

\medskip

Interpretation:
\begin{itemize}
\item \(S\) measures how parameter changes alter the normalized state,
\item \(S\) is the metric tensor on the variational manifold,
\item \(S\) is positive semidefinite.
\end{itemize}
\end{frame}

\begin{frame}{Deriving the generalized force \(g_i\)}
The right-hand side is
\[
\bra{\overline{\Psi}_i}(H-E)\ket{\overline{\Psi}}.
\]

Define the local energy
\[
E_{\mathrm{loc}}(x)=\frac{(H\Psi)(x)}{\Psi(x)}.
\]

Then
\[
g_i
=
\Re\left[
\bra{\overline{\Psi}_i}(H-E)\ket{\overline{\Psi}}
\right]
=
\langle O_i E_{\mathrm{loc}}\rangle
-
\langle O_i\rangle \langle E_{\mathrm{loc}}\rangle
\]
for real-valued states.

Thus \(g_i\) is the covariance between the logarithmic derivative and the local energy.
\end{frame}

\begin{frame}{Final SR linear system}
Putting everything together yields
\[
\sum_j S_{ij}\,\delta\alpha_j = -\delta\tau\, g_i,
\]
with
\[
S_{ij}
=
\langle O_i O_j\rangle
-
\langle O_i\rangle \langle O_j\rangle,
\]
and
\[
g_i
=
\langle O_i E_{\mathrm{loc}}\rangle
-
\langle O_i\rangle \langle E_{\mathrm{loc}}\rangle.
\]

In vector form:
\[
S\,\delta\bm{\alpha} = -\delta\tau\, \bm{g},
\qquad
\delta\bm{\alpha} = -\delta\tau\, S^{-1}\bm{g}.
\]
\end{frame}

\begin{frame}{Why this differs from ordinary gradient descent}
Ordinary gradient descent updates
\[
\delta\bm{\alpha}_{\mathrm{GD}} = -\eta \,\bm{g}.
\]

SR instead uses
\[
\delta\bm{\alpha}_{\mathrm{SR}} = -\eta \, S^{-1}\bm{g}.
\]

The inverse metric \(S^{-1}\) rescales and mixes parameter directions according to the geometry of the wave-function manifold.

\medskip

Consequences:
\begin{itemize}
\item reduced sensitivity to redundant parameterizations,
\item better conditioning,
\item close relation to imaginary-time evolution.
\end{itemize}
\end{frame}

\begin{frame}{Connection to the natural gradient}
The squared distance between nearby normalized states is
\[
ds^2
=
\sum_{ij} S_{ij}\, d\alpha_i\, d\alpha_j.
\]

Thus \(S\) defines a Riemannian metric in parameter space.

\medskip

Steepest descent in this metric gives
\[
\delta\bm{\alpha} \propto -S^{-1}\bm{g},
\]
which is exactly the \textbf{natural gradient} update.

Therefore SR is the natural-gradient method specialized to variational quantum states.
\end{frame}

\begin{frame}{Regularization and practical solution}
In practice \(S\) may be singular or ill-conditioned.

A standard fix is diagonal regularization:
\[
S \;\to\; S + \lambda I,
\qquad \lambda>0.
\]

Then solve
\[
(S+\lambda I)\,\delta\bm{\alpha} = -\delta\tau\,\bm{g}.
\]

Interpretation:
\begin{itemize}
\item stabilizes inversion,
\item damps noisy directions,
\item interpolates between SR and ordinary gradient descent for large \(\lambda\).
\end{itemize}
\end{frame}

\begin{frame}{Algorithmic summary}
At each iteration:
\begin{enumerate}
\item Sample configurations \(x\) from \(P(x)=|\Psi(x)|^2/\braket{\Psi}{\Psi}\).
\item Compute logarithmic derivatives \(O_k(x)\).
\item Compute local energies \(E_{\mathrm{loc}}(x)\).
\item Estimate
\[
S_{ij}=\langle O_i O_j\rangle - \langle O_i\rangle\langle O_j\rangle,
\]
\[
g_i=\langle O_i E_{\mathrm{loc}}\rangle - \langle O_i\rangle\langle E_{\mathrm{loc}}\rangle.
\]
\item Solve
\[
(S+\lambda I)\delta\bm{\alpha}=-\delta\tau\,\bm{g}.
\]
\item Update \(\bm{\alpha}\leftarrow \bm{\alpha}+\delta\bm{\alpha}\).
\end{enumerate}
\end{frame}

\begin{frame}{Summary of the derivation}
\begin{itemize}
\item Imaginary-time evolution suggests the optimal direction toward the ground state.
\item The variational manifold restricts admissible states to \(\ket{\Psi(\bm{\alpha})}\).
\item Projecting the infinitesimal imaginary-time step onto the tangent space yields
\[
S\,\delta\bm{\alpha}=-\delta\tau\,\bm{g}.
\]
\item \(S\) is the covariance matrix of logarithmic derivatives.
\item \(\bm{g}\) is the covariance of logarithmic derivatives with the local energy.
\item SR is both:
\begin{itemize}
\item a projected imaginary-time evolution,
\item and a natural-gradient optimization method.
\end{itemize}
\end{itemize}
\end{frame}

\end{document}
